\documentclass{article}
\usepackage{graphicx}

\usepackage[utf8]{inputenc}     % pdfLaTeX
\usepackage[T1]{fontenc}
\usepackage[magyar]{babel}

\title{Érzékszervek}
\author{Haydu Lőrinc, Lengyel Zoltán Péter, Rovenszky Robin}
\date{Legutóbb frissítve: 2026. Április 15.}

\begin{document}

\maketitle
\tableofcontents
\newpage

\section{Bevezető}
Az érszékszerveink receptorokon keresztül érzékelik az ún. adekvált ingert.

\begin{itemize}
    \item Photoreceptor - Szem
    \item Chemoreceptor - Szaglás, ízlelés
    \item Mechanoreceptor - Hallás, tapintás
\end{itemize}

\section{A látás}
\subsection{Szem}

% insert szem figures here

\begin{itemize}
    \item Pupilla $\rightarrow$ Szivárványhártya $\rightarrow$ Ínhártya
    \item Pupillareflex\\ A pupillánk szűkül világosabb, illetve tágul sötétebb helyzetekben, a retina védelmének érdekében.
\end{itemize}

% finish table TBA
\begin{tabular}{|l|l|l|} \hline
     & Pálcika & Csap \\ \hline
     Méret & hosszabb & \\
\end{tabular}
\

% színtévesztés, színvakság, zöldhólyag, szürkehólyag

\subsubsection{Színtévesztés}
\begin{itemize}
    \item A leggyakoribb a vörös-zöld színtévesztés (férfiakban genetikailag még gyakoribb)
    \item insert other stuff here
\end{itemize}

\subsubsection{Érzéki csalódások}
\begin{itemize}
    \item Az agyunk kiegészíti amiz objektíven látunk
    \item A mozgás illúziója: Zoetróp
\end{itemize}

\section{Hallás}
\subsection{A fül}

% insert fül figures here

\begin{itemize}
    \item A hallószerv három részre osztható: 
    \begin{itemize}
        \item Külső fül \begin{itemize}
            \item Fülkagyló
            \item Külső hallójárat
        \end{itemize}
        \item Középfül \begin{itemize}
            \item Dobhártya 
            \item 3 hallócsont \begin{itemize}
                \item Kalapács (dobhártyához kapcsolódik)
                \item Üllő 
                \item Kengyel
            \end{itemize}
        \end{itemize}
        \item Belső fül \begin{itemize}
            \item Csiga - folyadékkal tültütt üregrendszer \begin{itemize}
                \item Ovális ablak
                \item Felső járat
                \item Középső járat -  belsejében mechanoreceptorok (Corti-féle szerv: érzékszőrök + fedő)
                \item Alsó járat
                \item A hang elindul a felső járaton, a csiga végén visszafordul és az alsó járaton jön vissza
                \item Kerek ablak
            \end{itemize}
            \item Fülkürt (garattal van összeköttetésben, dobhártya két oldala közti légnyomást egyenlíti ki - nyelés)
        \end{itemize}
    \end{itemize}

    \item 400Hz-22000Hz között érzékel \begin{itemize}
        \item Mély hangok - a csiga csúcsa környékén váltanak ki ingerületet
        \item Magas hangok -  ovális ablakhoz közel váltanak ki ingerületet
    \end{itemize}
\end{itemize}

%dB table here

\subsection{Egyensúly érzékelő szervrendszer}
\begin{itemize}
    \item Mechanoreceptorok \begin{itemize}
        \item Ampulla - ívjárat alapi része közelében \begin{itemize}
            \item Szőrsejtek 
            \item Kupula - szőrsejteket körbevevő kocsonyás kúp
        \end{itemize}
    \end{itemize}
    \item Tömlőcske és zsákocska között \begin{itemize}
        \item Nagyobb, szélesebb alapon elhelyezkedő sejtek, tetején mészkristályok nyomják a szőrsejteket
    \end{itemize}
\end{itemize}

\section{Ízérzékelés}
\begin{itemize}
    \item Chemoreceptorok
    \item Sós - nátrium (pl. nátrium klorid, nátrium cegyületek)
    \item Édes - mono és diszacharidok 
    \item Savanyú - savak
    \item Keserű - általában szervezetet károsító anyagok
    \item Umami - nukleotidok, aminosavak közül glutaminsav (nátrium-glutamát ízfokozó)
    \item Újabb kutatások szerint vannak vizet érzékelő receptorok is
\end{itemize} 

\section{Szaglás}
\begin{itemize}
    \item Chemoreceptorok
    \item Légnemű anyagok érzékelése alkalmas
    \item Szaglóhám az orrüregben 
    \item A szaglóinger fáradásra hajlamos, ahhoz hogy fennmaradjon az inger, folyamatosan növekvő koncentráció szükséges (avagy erősödnie kell a szagnak, mert egy idő után folyamatos egyenletes szagot nem érzünk)
\end{itemize}

\section{Tapintás}

\end{document}